\documentclass[a4paper,12pt]{article}
\usepackage[utf8]{inputenc}
\usepackage[T1]{fontenc}
\usepackage[ngerman]{babel}
\usepackage{amsmath}
\usepackage{amssymb}
\usepackage{enumitem}
\usepackage{geometry}
\geometry{a4paper, margin=2.5cm}

\title{Einsendeaufgaben -- Lektion 2\\[0.5em]
\large Modul 61112: Lineare Algebra}
\author{}
\date{}

\begin{document}
\maketitle

\section*{Aufgabe 2.3}

Es existiert wegen Korollar 2.2.13 eine geordnete Basis $B_U = (v_1, \ldots, v_t)$ für $U$.
Nach dem Basisergänzungssatz lässt sich $B_U$ erweitern zu einer geordneten Basis
$\mathcal{B} = (v_1, \ldots, v_t, v_{t+1}, \ldots, v_n)$ für $V$.

Wir definieren die Linearformen $a_1, \ldots, a_{n-t} \in V^*$ durch die Basis $\mathcal{B}$.
Denn eine Linearform ist eine lineare Abbildung und es reicht demnach aus,
die Abbildung der Basis zu bestimmen.

Also definieren wir $a_i$ mit $i \leq n-t$ wie folgt:
\[
  a_i(v) = \begin{cases} 1, & \text{wenn } v = v_{t+i} \\ 0, & \text{sonst} \end{cases}
\]

Die Basis von $\ker(a_i)$ für ein $i \leq n-t$ ist nach Definition $\mathcal{B} \setminus \{v_{t+i}\}$.
Demnach gilt
\[
  \bigcap_{i=1}^{n-t} \mathcal{B} \setminus \{v_{t+i}\} = B_U \cap \bigcap_{i=1}^{n-t} \{v_{t+i}\} = B_U.
\]

Daraus folgt:
\[
  \bigcap_{i=1}^{n-t} \ker(a_i)
  = \bigcap_{i=1}^{n-t} \langle \mathcal{B} \setminus \{v_{t+i}\} \rangle
  = \langle B_U \rangle
  = U
\]

\end{document}
